\section{量子化复流形}\label{qcomp}

本节我们转向一个多少有些不同的方向。设 $Y$ 是复辛流形，带全纯辛形式 $\Omega$。通过忘记其复结构，可把 $Y$ 视为实辛流形。我们想用量子化来解释复流形 $Y$ 的量子化，即把 $Y$ 的量子化定义为把 $Y$ 视为实流形时的膜量子化。

沿用第 \ref{amo} 节的记号，$Y$ 作为实辛流形的典范复化是 $\h Y=Y_1\times Y_2$，其中 $Y_1$ 与 $Y_2$ 都同构于 $Y$，但复结构相反。因此 $\h Y$ 作为实流形带有两个复结构 $I_1=I\oplus I$ 与 $I_2=-I\oplus(-I)$。交换两个因子的对合
\be\label{qcomplextwo}\tau:(y_1,y_2)\mapsto(\bar y_2,\bar y_1)\ee
是反全纯的，其不动点集是对角 $\Delta\cong Y$。这是 $\h Y$ 中实维数为 $\dim_\C Y$ 的实子流形。$\tau$ 与 $\h Y$ 上的复辛形式
\be\label{qcomplexone}\h\Omega=\frac12\Omega\boxplus\frac12\bar\Omega\ee
相容，其中 $\Omega$ 是 $Y_1$ 上的全纯辛形式，$\bar\Omega$ 是 $Y_2$ 上的全纯辛形式（在 $Y_2$ 的复结构下）。限制到对角 $\Delta$ 上，$\h\Omega$ 等于 $\Re\Omega$。因此我们取 $\h Y$ 的实辛形式为 $\Im\h\Omega$，使 $\Delta$ 为拉格朗日。相应地，B-场为 $\Re\h\Omega=\frac12\Re\Omega\boxplus\frac12\Re\Omega$，限制到 $Y$ 上就是 $\Re\Omega$。

现在我们来描述 $\h Y$ 上的典范余各向同性膜 $\h\B_\cc$。在 $Y_1$ 上，$\sF+\sB=\frac12\Re\Omega=\frac12\omega_J$，复合结构 $\omega^{-1}(\sF+\sB)$ 等于 $(\frac12\omega_K)^{-1}(\frac12\omega_J)=I$，其中 $I$ 是 $Y_1$ 的复结构。在 $Y_2$ 上，复结构是 $-I$，而 $\omega^{-1}(\sF+\sB)=(\frac12\omega_K)^{-1}(\frac12\omega_J)=I$，故相对于 $Y_2$ 的复结构，它等于 $-I$。因此在 $Y_1$ 与 $Y_2$ 上，$\sF+\sB$ 都给出可积复结构，$\B_{\cc,1}$ 与 $\B_{\cc,2}$ 都是 $\h Y$ 上的典范余各向同性膜，其 $\CP$ 丛平凡。乘积 $\h\B_\cc=\B_{\cc,1}\times\B_{\cc,2}$ 是 $\h Y$ 上的典范余各向同性膜。

$\h Y$ 上的拉格朗日膜 $\h\B$ 支撑在对角 $\Delta\cong Y$ 上。为满足 $\sF+\sB=0$，需取 $\sF=-\Re\Omega$。由于 $\frL$ 是 $Y$ 上曲率为 $\Re\Omega$ 的预量子线丛，$\h\B$ 的 $\CP$ 丛可取为 $\frL^{-1}$。与通常一样，反常条件因复辛流形有典则 spin 结构而自动满足。于是所预想的量子 Hilbert 空间是
$\H=\Hom(\h\B,\h\B_\cc)$。

与形变量子化的关系是清楚的。我们有
\[\Hom(\h\B_\cc,\h\B_\cc)=\Hom(\B_{\cc,1},\B_{\cc,1})\times\Hom(\B_{\cc,2},\B_{\cc,2}).\]由第 \ref{abranes} 节，$\Hom(\B_{\cc,1},\B_{\cc,1})=\A$ 是 $Y$ 上全纯函数的形变量子化，而 $\Hom(\B_{\cc,2},\B_{\cc,2})=\bar\A$ 是反全纯函数的形变量子化。因此 $\Hom(\h\B_\cc,\h\B_\cc)=\A\times\bar\A$，正如把 $Y$ 视为实辛流形时 $\Hom$ 代数所应为。

我们可以像 §\ref{abranes} 那样展开这个构造。假设 $\frL=\L^2$ 是某个线丛的平方，使 $c_1(\frL)$ 为偶。作一个 B-场变换，可令 $\sB=0$，于是两个 $\B_\cc$ 的 $\CP$ 丛都是 $\L$（其曲率为 $\frac12\omega_J$），而 $\h\B$ 的 $\CP$ 丛平凡。$(\h\B,\h\B_\cc)$ 弦的世界面是条带 $\R\times[0,\pi]$：在 $\sigma=0$ 处边界条件由 $\h\B_\cc$ 给出，它是乘积膜，与展开无关；在 $\sigma=\pi$ 处边界条件由 $\h\B$ 给出，它不是乘积膜，把 $Y$ 的两个拷贝粘在一起。

因为 $\h\B$ 支撑在对角上，展开后得到 $[0,2\pi]$ 上单个 $Y$ 的 $\sigma$-模型，其边界条件在 $\sigma=0$ 与 $\sigma=2\pi$ 处由两个（不同的）膜给出。具体地，把 $\sigma>\pi$ 的部分翻转，会改变其定向，等价于改变辛形式的符号；这使两半上的辛形式都变为 $\omega_K$。边界定向的反转把 $\CP$ 丛取逆，故在 $\sigma=2\pi$ 处膜的 $\CP$ 曲率为 $-\frac12\omega_J$，即它是共轭膜 $\bar\B_\cc$。于是展开后的 $\Hom(\h\B,\h\B_\cc)$ 被重新解释为 $\Hom(\bar\B_\cc,\B_\cc)$。

我们本可以先取 $\sB\ne0$ 再展开，但那样 B-场在条带中央会不连续，该不连续处成为一个拓扑线算符。先令 $\sB=0$ 的代价是假设 $\frL$ 有平方根。在与几何 Langlands 相关的情形（\cite{GaW}），$\omega_J$ 的上同调平凡，而 $\frL$ 平凡，故无 B-场的版本更为方便。

\begin{yizhu}[本节与几何 Langlands 的关系]
把复辛流形 $Y$ 当作“实辛流形 $Y$”来量子化，听起来玄，其实是几何 Langlands 里的标准图像：Hitchin 模空间既是复辛流形，又是某个实辛流形的复化；其 $A$-模型给出 $D$-模侧，$B$-模型（全纯函数的形变量子化）给出 Higgs 丛侧。原文 §\ref{flatg} 中“$H^0(M,K_M^{1/2}\otimes\frL)$ 与复结构 $t$ 无关”的陈述，正是这一图像最可检验的推论之一；沿 Teichmüller 空间移动 $t$ 给出的平坦联络（ADW、Hitchin、Faltings），则是它的几何实现。[来源：A 级]
\end{yizhu}

\medskip
\noindent{\bfseries 致谢}\quad
我们感谢 M. Abouzeid、R. Bezrukavnikov、K. Costello、E. Frenkel、P. Etingof、V. Ginzburg、D. Kazhdan、S. Li 与 R. Mazzeo 的讨论。

Perimeter Institute 的研究由加拿大政府通过 Industry Canada 以及安大略省研究与创新部资助。E. Witten 的研究部分由 NSF 基金 PHY-1911298 资助。
